\chapter{Related Work}
\label{ch:related_work}

This chapter situates the thesis within the existing literature. We review, in turn, deep learning for chest radiography and the public benchmarks that enabled it, efficient and cascaded architectures, foundation models for medical imaging, uncertainty estimation and calibration, conformal prediction, synthetic data and augmentation, and explainability. The chapter closes by positioning the proposed system relative to this body of work.

\section{Deep Learning in Chest Radiography}
The modern era of medical image classification began with the success of deep convolutional neural networks on natural images, starting with AlexNet \cite{krizhevsky2012imagenet} and continuing through VGG \cite{simonyan2015very}, residual networks \cite{he2016deep}, and densely connected networks \cite{huang2017densely}. The availability of very large labelled chest-radiograph datasets made it possible to transfer these architectures to the medical domain. The NIH ChestX-Ray14 dataset \cite{wang2017chestxray}, with over one hundred thousand frontal radiographs labelled for fourteen pathologies, and the Stanford CheXpert dataset \cite{irvin2019chexpert}, with its explicit treatment of label uncertainty, became the standard benchmarks; MIMIC-CXR \cite{johnson2019mimic} later added free-text reports at a similar scale.

Building on these resources, CheXNet \cite{rajpurkar2017chexnet} demonstrated radiologist-level pneumonia detection with a 121-layer DenseNet, and a large number of follow-up studies refined the architecture, the loss function, and the labelling strategy. Almost all of these systems share two characteristics that this thesis sets out to change: they are single-stage, applying one fixed network to every image, and they report only a point prediction, with no calibrated uncertainty and no mechanism to abstain. Broad surveys of deep learning in medical image analysis \cite{litjens2017survey} and landmark results in other modalities, such as dermatologist-level skin-cancer classification \cite{esteva2017dermatologist}, confirm both the promise of the approach and the recurring concern that such models are poorly calibrated and difficult to trust in isolation.

\section{Efficient and Cascaded Architectures}
A parallel line of research has sought to reduce the computational cost of deep networks. MobileNets \cite{howard2017mobilenets, sandler2018mobilenetv2} use depthwise-separable convolutions and inverted residual blocks to achieve a strong accuracy-to-cost ratio on mobile and embedded hardware, while EfficientNet \cite{tan2019efficientnet} introduces compound scaling to balance network depth, width, and input resolution. More recently, vision transformers \cite{dosovitskiy2021image} and their hierarchical variant, the Swin Transformer \cite{liu2021swin}, have matched or exceeded convolutional networks at the cost of larger models and more data.

The idea of routing inputs through models of increasing capacity --- variously called cascades, early-exit networks, or conditional computation --- predates deep learning but fits it naturally: most inputs are easy and can be resolved cheaply, so spending the full model only on the hard cases reduces the average cost. The present work adopts exactly this principle, pairing a MobileNetV2 screening stage with a heavier EfficientNet-B4 or Swin-based stage, and adds a confidence-based router that decides, per image, whether escalation is warranted.

\section{Foundation Models for Medical Imaging}
Transfer learning from ImageNet \cite{deng2009imagenet, russakovsky2015imagenet} has long been the default initialisation for medical image classifiers, but the domain gap between natural and medical images limits its benefit. This has motivated medical-specific pretraining. The Ark family of models \cite{ma2023foundation} accrues and reuses knowledge across many public chest-radiograph datasets and their heterogeneous label sets, producing a backbone that transfers more robustly to downstream chest tasks than a generically pretrained network. In this thesis, the Ark+ model with a Swin Transformer backbone is used as one of the two Tier-2 options, and is compared head-to-head with a conventionally pretrained EfficientNet-B4.

\section{Uncertainty Estimation and Calibration}
Quantifying predictive uncertainty is essential for safe deployment. Gal and Ghahramani \cite{gal2016dropout} showed that dropout, when left active at test time, can be interpreted as approximate Bayesian inference: performing $T$ stochastic forward passes (Monte Carlo Dropout) yields a distribution over predictions whose variance estimates epistemic uncertainty. Deep ensembles \cite{lakshminarayanan2017simple} provide a strong, if expensive, alternative, and Kendall and Gal \cite{kendall2017uncertainties} formalise the distinction between aleatoric and epistemic uncertainty in deep vision models. Test-time augmentation, which aggregates predictions over augmented copies of an input, is a complementary technique that also exposes the sensitivity of a prediction to plausible input perturbations \cite{wang2019aleatoric}.

A related but distinct problem is calibration: whether a model's stated confidence matches its empirical accuracy. Modern deep networks are known to be poorly calibrated \cite{guo2017calibration}, typically over-confident, and temperature scaling is a simple and effective post-hoc remedy. Calibration quality is commonly summarised by the Expected Calibration Error \cite{naeini2015obtaining}. This thesis applies Monte Carlo Dropout and test-time augmentation for uncertainty, temperature scaling for calibration, and reports the Expected Calibration Error throughout.

\section{Conformal Prediction}
Conformal Prediction \cite{vovk2005algorithmic, shafer2008tutorial} is a framework that turns any point predictor into a set predictor with a finite-sample, distribution-free coverage guarantee: for a chosen error rate $\alpha$, the produced prediction set contains the true label with probability at least $1-\alpha$, under only the assumption of exchangeability. In classification, this means the model returns a set of plausible labels rather than a single forced decision, and the size of that set is itself an interpretable measure of difficulty. Recent work has made the framework practical and adaptive for deep classifiers, including the adaptive prediction sets of Romano et al.\ \cite{romano2020classification} and the regularised image-classifier sets of Angelopoulos et al.\ \cite{angelopoulos2021uncertainty}; a comprehensive and accessible treatment is given by Angelopoulos and Bates \cite{angelopoulos2023gentle}. The present system calibrates a conformal predictor on a held-out split and uses the resulting prediction sets both as an output to the clinician and as a signal for abstention.

\section{Synthetic Data and Augmentation}
Class imbalance is endemic in medical imaging: pathological findings are, by definition, rarer than normal studies. Data augmentation is the standard first response, and surveys \cite{shorten2019survey} catalogue the geometric and photometric transforms commonly used. Stronger regularisers that mix examples, such as mixup \cite{zhang2018mixup} and CutMix \cite{yun2019cutmix}, improve both accuracy and calibration and are evaluated in this thesis. Generative models offer a more aggressive route to balancing the minority class: generative adversarial networks were used early for medical image synthesis \cite{goodfellow2014generative, fridadar2018gan}, and diffusion models --- denoising diffusion probabilistic models \cite{ho2020denoising} and the latent-diffusion formulation behind Stable Diffusion \cite{rombach2022high} --- now produce high-fidelity samples. We use Stable Diffusion to synthesise additional pneumothorax-positive images, gated by a Fr\'echet Inception Distance quality threshold, and quantify the effect with dedicated ablations.

\section{Explainability}
For a clinical decision-support tool, a prediction is more useful when it is accompanied by a visual explanation. Class Activation Mapping \cite{zhou2016learning} and its gradient-based generalisation Grad-CAM \cite{selvaraju2017gradcam} highlight the image regions most responsible for a prediction. HiResCAM \cite{draelos2020hirescam} is a more recent variant that is provably faithful to the model for certain architectures, avoiding some of the spatial artefacts of Grad-CAM. This thesis generates HiResCAM saliency maps for the active model in each prediction.

\section{Positioning of This Work}
The literature contains strong individual components: efficient backbones, medical foundation models, Monte Carlo Dropout and ensembles, calibration methods, conformal prediction, diffusion-based augmentation, and faithful saliency. What is comparatively rare is a single, fully integrated system that combines all of them around an adaptive routing mechanism and then measures the contribution of each part. This thesis provides exactly such an integration, together with a genuine fifteen-configuration ablation study and a cross-dataset evaluation, in a deployable software package.
